\subsubsection{Sprint 9}

\subsubsubsection{Informazioni generali e attività da svolgere}

\begin{tabular}{p{0.25\textwidth} p{0.2\textwidth}}
	\textbf{Inizio} & 09/03/2026 \\
	\textbf{Fine prevista} & 16/03/2026 \\
	\textbf{Fine reale} & - \\
	\textbf{Giorni di ritardo} & -
\end{tabular}
\newline \newline 
Le nostre attività da svolgere, definite nel \href{https://nullpointersgroup.github.io/Documentazione/output/PB/Verbali%20Interni/2026-03-09_verbale_interno.pdf}{Verbale interno del 09/03/2026}, nel \href{https://nullpointersgroup.github.io/Documentazione/output/PB/Verbali%20Interni/2026-03-11_verbale_interno.pdf}{Verbale interno del 11/03/2026} e nel \href{https://nullpointersgroup.github.io/Documentazione/output/PB/Verbali%20Esterni/2026-03-11_verbale_esterno.pdf}{Verbale esterno del 11/03/2026} sono:
\begin{itemize}[itemsep=5pt, parsep=1pt, label=$\scriptstyle\bullet$]
	\item Sistemare i riferimenti informativi e normativi
	\item Scrivere il consuntivo$^G$ dello sprint$^G$ 8 e il preventivo dello sprint$^G$ 9
	\item Stesura Diario di Bordo del 11 marzo 2026
	\item Modellare il diagramma ER del DB$^G$ relazionale
	\item Inserire termini Glossario
	\item Correggere i vincoli di requisito dell'Analisi dei Requisiti$^G$
\end{itemize}

\subsubsubsection{Rischi attesi}
I possibili rischi attesi sono:
\begin{itemize}[itemsep=5pt, parsep=1pt, label=$\scriptstyle\bullet$]
	\item RP3: rischio d'incomprensione di qualche task$^G$
	\item RA1: rischio di qualche imprevisto personale
	\item RE1: rischio dovuto alla poca conoscenza sulle nuove tecnologie da affrontare
	\item RE2: rischio dovuto alla stima errata delle attività
\end{itemize}

\subsubsubsection{Preventivo}
Si prospetta l'utilizzo delle seguenti risorse:
\begin{table}[H]
	\centering
	\makebox[\textwidth][c]{%
		\begin{minipage}{1.2\textwidth}
			\centering
			% Prima tabella senza bordi per le intestazioni ruotate
			\begin{tabular}{p{0.2\textwidth}p{0.12\textwidth}p{0.12\textwidth}p{0.12\textwidth}p{0.12\textwidth}p{0.12\textwidth}p{0.12\textwidth}}
				& \rotatebox{45}{\textbf{Responsabile$^G$}} &
				\rotatebox{45}{\textbf{Amministratore$^G$}} & \rotatebox{45}{\textbf{Analista$^G$}} & \rotatebox{45}{\textbf{Programmatore$^G$}} & \rotatebox{45}{\textbf{Verificatore$^G$}} &
				\rotatebox{45}{\textbf{Progettista$^G$}} \\
			\end{tabular}
			
			\vspace{0.2cm} % Piccolo spazio tra le due tabelle
			
			% Seconda tabella con i dati e i bordi
			\begin{tabular}{|p{0.2\textwidth}|p{0.12\textwidth}|p{0.12\textwidth}|p{0.12\textwidth}|p{0.12\textwidth}|p{0.12\textwidth}|p{0.12\textwidth}|}
				\hline
				M. Mazzaretto & 3 & - & - & - & - & - \\ \hline
				T. Ceron      & - & - & - & 3 & - & 3 \\ \hline
				L. Pieripolli & - & - & - & - & - & 3 \\ \hline
				L. Marcuzzo   & - & - & - & - & - & 3 \\ \hline
				M. Brunello   & - & - & - & - & 3 & 3 \\ \hline
				L. Casagrande & - & 1 & - & - & - & 3 \\ \hline
			\end{tabular}
		\end{minipage}
	}
	\caption{Sprint 9: Preventivo}
\end{table}

Il totale preventivato per lo sprint$^G$ è \textbf{480€}.\\

\begin{figure}[H]
	\centering
	\includegraphics[width=0.8\textwidth]{PianoProgetto/sprint09_preventivo}
	\caption{Sprint 9: Preventivo}
\end{figure}